\begin{abstract}
This paper evaluates a classical single-channel EEG pipeline for two related sleep tasks: sleep-stage classification and apnea detection. The study is centered on three traditional model families, namely SVM with RBF kernel, Random Forest, and XGBoost, under strictly subject-wise validation and cross-dataset transfer protocols. For sleep staging, the main in-dataset reference is \expname{sleep_edf_expanded_tuned}, while external generalization is measured with \expname{cross_sleep_edf_to_isruc_n123} and \expname{cross_isruc_to_sleep_edf_n123}. For apnea detection, the main multitarget reference is \expname{mitbih_apnea_stage_classic}, complemented with the cross-dataset experiments \expname{cross_dataset_mitbih_to_st_vincent_classic} and \expname{cross_dataset_st_vincent_to_mitbih_classic}. The pipeline uses 30-second windows, shared tabular EEG features, and metrics aligned with the course brief: accuracy, macro-F1, and Cohen's kappa for staging; accuracy, sensitivity, specificity, and AUC-ROC for apnea. The main findings are the following. First, XGBoost achieved the best in-dataset staging accuracy and kappa on Sleep-EDF, while Random Forest achieved the best macro-F1. Second, XGBoost also offered the best overall balance on MIT-BIH when apnea and staging were considered jointly, although paired AUC analysis showed that the apnea advantage over Random Forest was not robust when pooled out-of-fold predictions were compared. Third, all cross-dataset settings showed a pronounced performance drop, indicating fragile generalization under domain shift. These results support the use of classical monocanal EEG baselines as reproducible references, but they also show that domain adaptation is necessary before claiming robust deployment across cohorts.
\end{abstract}

\begin{IEEEkeywords}
sleep staging, sleep apnea, EEG, single-channel EEG, subject-wise validation, cross-dataset generalization, random forest, xgboost, support vector machine
\end{IEEEkeywords}

\section*{Resumen}
Este trabajo eval\'ua una l\'inea cl\'asica de EEG monocanal para dos tareas relacionadas del sue\~no: clasificaci\'on de etapas y detecci\'on de apnea. La evaluaci\'on se centra en tres familias de modelos tradicionales: \emph{SVM-RBF}, \emph{Random Forest} y \emph{XGBoost}, bajo protocolos \emph{subject-wise} estrictos y experimentos \emph{cross-dataset}. Para staging, el baseline intra-dataset principal es \expname{sleep_edf_expanded_tuned}, mientras que la generalizaci\'on externa se mide con \expname{cross_sleep_edf_to_isruc_n123} y \expname{cross_isruc_to_sleep_edf_n123}. Para apnea, el baseline multitarget principal es \expname{mitbih_apnea_stage_classic}, complementado con \expname{cross_dataset_mitbih_to_st_vincent_classic} y \expname{cross_dataset_st_vincent_to_mitbih_classic}. El pipeline usa ventanas de 30 s, \feature{eeg\_*} compartidas y m\'etricas alineadas con el laboratorio: \emph{accuracy}, \emph{macro-F1} y \emph{kappa} para staging; \emph{accuracy}, sensibilidad, especificidad y AUC-ROC para apnea. Los hallazgos principales fueron tres. Primero, XGBoost obtuvo el mejor \emph{accuracy}/\emph{kappa} intra-dataset en Sleep-EDF, mientras que Random Forest obtuvo el mejor \emph{macro-F1}. Segundo, en MIT-BIH, XGBoost mostr\'o el mejor balance general entre apnea y staging, aunque el contraste pareado sobre AUC sugiri\'o que la ventaja frente a Random Forest no es estable bajo todas las formas de agregaci\'on. Tercero, todos los escenarios \emph{cross-dataset} sufrieron una ca\'ida marcada del rendimiento, lo que evidencia fragilidad ante cambio de dominio. En conjunto, los resultados validan la utilidad de los baselines cl\'asicos sobre EEG monocanal como referencia reproducible, pero tambi\'en justifican una mejora metodol\'ogica orientada a adaptaci\'on de dominio.
